\begin{figure}[t]
\centering

% --------------------------------------------------
% Shared styles
% --------------------------------------------------
\tikzset{
    roundingfig/.style={
        >=Latex,
        line cap=round,
        line join=round,
        every node/.style={font=\small},
        lattice/.style={circle, fill=black!30, inner sep=1.05pt},
        feasiblept/.style={circle, fill=black, inner sep=1.45pt},
        lpregion/.style={draw=none, fill=blue!8},
        constraint/.style={blue!70!black, line width=1.0pt},
        solLP/.style={circle, fill=blue!80!black, inner sep=2pt},
        roundpt/.style={circle, fill=red!70!black, inner sep=2pt},
        pertpt/.style={circle, fill=brown!70!black, inner sep=2pt},
        fppt/.style={circle, fill=green!55!black, inner sep=2pt},
        projpt/.style={circle, fill=purple!80!black, inner sep=2pt}
    }
}

% Shared zoom window
\def\zxmin{2.5}
\def\zxmax{4.5}
\def\zymin{1.75}
\def\zymax{3.25}

% --------------------------------------------------
% Macro: shared geometry
% --------------------------------------------------
\newcommand{\DefineRoundingGeometry}{%
    \def\eps{0.20}%
    \pgfmathsetmacro{\aone}{(-1 - 2*\eps)/4}%
    \pgfmathsetmacro{\bone}{3 + \eps}%
    \pgfmathsetmacro{\atwo}{(-3 + 2*\eps)/2}%
    \pgfmathsetmacro{\btwo}{-\atwo*(5+\eps)}%
    \pgfmathsetmacro{\xLP}{(\btwo-\bone)/(\aone-\atwo)}%
    \pgfmathsetmacro{\yLP}{\aone*\xLP+\bone}%
    \pgfmathsetmacro{\xD}{-\btwo/\atwo}%
    \pgfmathsetmacro{\xproj}{2.78}
    \pgfmathsetmacro{\yproj}{\aone*\xproj + \bone}
    % Chosen illustrative points
    \coordinate (xbar)   at (\xLP,\yLP);%
    \coordinate (xround) at (4,2);%
    \coordinate (xpert)  at (3,3);
    \coordinate (xproj)  at (\xproj,\yproj);%
    \coordinate (xfp)    at (3,2);%
}

% --------------------------------------------------
% Macro: draw feasible region and lattice
% --------------------------------------------------
\newcommand{\DrawRoundingBase}{%
    % LP relaxation
    \filldraw[lpregion]
        (0,0) --
        (\xD,0) --
        (\xLP,\yLP) --
        (0,\bone) --
        cycle;

    % Constraint boundaries
    \draw[constraint] (0,\bone) -- (\xLP,\yLP);
    \draw[constraint] (\xLP,\yLP) -- (\xD,0);
    \draw[constraint] (0,0) -- (\xD,0);
    \draw[constraint] (0,0) -- (0,\bone);

    % Integer lattice points
    \foreach \x in {0,...,5}{%
      \foreach \y in {0,...,4}{%
        \pgfmathtruncatemacro{\ok}{%
          ifthenelse(%
            \y <= \aone*\x+\bone + 0.001 &&%
            \y <= \atwo*\x+\btwo + 0.001,%
            1,0)%
        }%
        \ifnum\ok=1
          \node[feasiblept] at (\x,\y) {};%
        \else
          \node[lattice] at (\x,\y) {};%
        \fi
      }%
    }%
}

% --------------------------------------------------
% Macro: draw zoom axes and ticks
% --------------------------------------------------
\newcommand{\DrawZoomAxes}{%
    \draw[->, thick]
      (\zxmin,\zymin) -- (\zxmax+0.22,\zymin) node[below right] {$x_1$};
    \draw[->, thick]
      (\zxmin,\zymin) -- (\zxmin,\zymax+0.22) node[above left] {$x_2$};

    \foreach \x/\lab in {3/3,4/4}{%
      \draw (\x,\zymin) -- (\x,\zymin-0.035) node[below] {\scriptsize \lab};%
    }%
    \foreach \y/\lab in {2/2}{%
      \draw (\zxmin,\y) -- (\zxmin-0.035,\y) node[left] {\scriptsize \lab};%
    }%
}

% ==================================================
% Global view
% ==================================================
\begin{subfigure}[t]{0.315\textwidth}
\centering
\vspace{0pt}%
\begin{tikzpicture}[roundingfig, x=0.8cm, y=0.8cm]

\DefineRoundingGeometry

% Fixed bounding box keeps the panel width predictable without scaling text.
\path[use as bounding box] (-0.35,-0.38) rectangle (6.40,4.65);

% Axes
\draw[->, thick] (-0.25,0) -- (5.75,0) node[below right] {$x_1$};
\draw[->, thick] (0,-0.25) -- (0,4.35) node[above left] {$x_2$};

\DrawRoundingBase

% Highlight zoom window
\draw[black!45, dashed, rounded corners=1pt]
    (\zxmin,\zymin) rectangle (\zxmax,\zymax);

% Points
\node[solLP] at (xbar) {};
\node[blue!80!black, below] at (xbar) {$\bar{x}$};

\node[roundpt] at (xround) {};
\node[red!70!black, above right] at (xround) {$x^{\mathrm{round}}$};

\node[fppt] at (xfp) {};
\node[green!45!black, below left] at (xfp) {$x^{\mathrm{FP}}$};

\end{tikzpicture}
\label{fig:rounding-vs-fp-global}
\end{subfigure}%
\hfill
% ==================================================
% Zoom: rounding
% ==================================================
\begin{subfigure}[t]{0.315\textwidth}
\centering
\vspace{0pt}%
\begin{tikzpicture}[roundingfig, x=2.0cm, y=2.0cm]

\DefineRoundingGeometry

% Fixed bounding box keeps both zoom panels identical.
\path[use as bounding box]
  (\zxmin-0.24,\zymin-0.27) rectangle (\zxmax+0.45,\zymax+0.36);

\DrawZoomAxes

\begin{scope}
  \clip (\zxmin,\zymin) rectangle (\zxmax,\zymax);

  \DrawRoundingBase

  \node[solLP] at (xbar) {};
  \node[blue!80!black, left, xshift=-2.5mm] at (xbar) {$\bar{x}$};

  \node[roundpt] at (xround) {};
  \node[red!70!black, above] at (xround) {$x^{\mathrm{round}}$};

  \draw[red!70!black, dashed, ->] (xbar) -- (xround);
\end{scope}

\end{tikzpicture}
\label{fig:rounding-zoom}
\end{subfigure}%
\hfill
% ==================================================
% Zoom: feasibility pump
% ==================================================
\begin{subfigure}[t]{0.315\textwidth}
\centering
\vspace{0pt}%
\begin{tikzpicture}[roundingfig, x=2.0cm, y=2.0cm]

\DefineRoundingGeometry

% Fixed bounding box keeps both zoom panels identical.
\path[use as bounding box]
  (\zxmin-0.24,\zymin-0.27) rectangle (\zxmax+0.45,\zymax+0.36);

\DrawZoomAxes

\begin{scope}
  \clip (\zxmin,\zymin) rectangle (\zxmax,\zymax);

  \DrawRoundingBase

  \node[solLP] at (xbar) {};
  \node[blue!80!black, left, xshift=-2.5mm] at (xbar) {$\bar{x}$};

  \node[roundpt] at (xround) {};
  \node[red!70!black, above right] at (xround) {$x^{I}$};

  \node[pertpt] at (xpert) {};
  \node[brown!80!black, above] at (xpert) {$x^{\mathrm{pert}}$};

  \node[projpt] at (xproj) {};
  \node[purple!80!black, above, yshift=2mm] at (xproj) {$x^{\mathrm{proj}}$};

  \node[fppt] at (xfp) {};
  \node[green!45!black, below left] at (xfp) {$x^{\mathrm{FP}}$};

  \draw[red!70!black, dashed, ->] (xbar) -- (xround);
  \draw[brown!80!black, dashed, ->] (xround) to (xpert);
  \draw[purple!80!black, dashed, ->] (xpert) to (xproj);
  \draw[red!55!black, dashed, ->] (xproj) to (xfp);
\end{scope}

\end{tikzpicture}
\label{fig:fp-zoom}
\end{subfigure}

\caption{Comparison between simple rounding and a schematic feasibility-pump trajectory.
Simple rounding may map the fractional LP solution $\bar{x}$ to an infeasible integer
point, whereas the feasibility pump alternates between integer rounding, perturbation,
and projection back onto the LP relaxation.}
\label{fig:rounding_vs_fp}
\end{figure}